% Generated from locale/tr/content/methods/induction/structural-induction.tex; do not hand-edit.


\olfileid[tr]{mth}{ind}{sti}

\olsection{Yapısal Tümevarım}

Şimdiye kadar tümevarımı bütün doğal sayılar hakkındaki sonuçları
belirlemek için kullandık. Fakat buna karşılık gelen bir ilke,
tümevarımsal olarak tanımlanmış bir kümenin bütün !!{element}s
hakkındaki sonuçları doğrudan kanıtlamak için kullanılabilir. Buna çoğu
kez \emph{yapısal} tümevarım denir; çünkü tümevarımsal olarak tanımlanan
nesnelerin yapısına dayanır.

Genellikle bir tümevarımsal tanım, (a)~kümenin ``başlangıçtaki''
!!{element}s listesinden ve (b)~eski !!{element}s arasından yenilerini üreten
işlemlerin bir listesinden oluşur. Örneğin düzgün terimler söz konusu
olduğunda başlangıç nesneleri harflerdir. Yalnızca bir işlemimiz vardır:
\begin{align*}
  o(s_1, s_2) = & [s_1 \circ s_2]
\end{align*}
Doğal sayıların~$\Nat$ kendilerinin bile tümevarımsal bir tanımla
verildiğini düşünebilirsiniz: başlangıç nesnesi~$0$, işlem ise ardıl
fonksiyonu~$x + 1$'dir.

Tümevarımsal olarak tanımlanmış bir kümenin bütün elemanları hakkında
bir şeyi, yani kümedeki her !!{element} için bir $P$ özelliğinin geçerli
olduğunu kanıtlamak için şunları yapmalıyız:
\begin{enumerate}
\item Başlangıç nesnelerinin~$P$'ye sahip olduğunu kanıtlamak.
\item Her~$o$ işlemi için, girdiler $P$'ye sahipse sonucun da sahip
  olduğunu kanıtlamak.
\end{enumerate}
Örneğin bütün düzgün terimler hakkında bir şeyi kanıtlamak için, onun
bütün harfler hakkında doğru olduğunu ve $s_1$ ile $s_2$ hakkında ayrı
ayrı doğru olması koşuluyla $[s_1 \circ s_2]$ hakkında da doğru olduğunu
kanıtlardık.

\begin{prop}
  Herhangi bir düzgün~$t$ terimindeki $[$ sayısı $]$ sayısına eşittir.
\end{prop}

\begin{proof}
Yapısal tümevarım kullanırız. Düzgün terimler, başlangıç nesneleri
harfler ve eski düzgün terimlerden yenilerini kuran işlem $o$ olmak
üzere, tümevarımsal olarak tanımlanmıştır.
\begin{enumerate}
\item İddia her harf için doğrudur; çünkü tek başına bir harfteki $[$
  sayısı~$0$, içindeki $]$ sayısı da~$0$'dır.
\item $s_1$'deki $[$ sayısının $]$ sayısına eşit olduğunu ve aynı şeyin
  $s_2$ için doğru olduğunu varsayın. $o(s_1, s_2)$'de, yani $[s_1
  \circ s_2]$'de bulunan $[$ sayısı, $s_1$ ile $s_2$'deki $[$
  sayılarının toplamına bir eklenmiş hâlidir. $o(s_1, s_2)$'deki $]$
  sayısı, $s_1$ ile $s_2$'deki $]$ sayılarının toplamına bir eklenmiş
  hâlidir. Dolayısıyla $o(s_1, s_2)$'deki $[$ sayısı, $o(s_1,s_2)$'deki
  $]$ sayısına eşittir.
\end{enumerate}
\end{proof}

\begin{prob}
  Hiçbir düzgün terimin~$]$ ile başlamadığını yapısal tümevarımla
  kanıtlayın.
\end{prob}

Yapısal tümevarımla başka bir kanıt verelim: bir simgeler dizisi~$t$'nin
öz başlangıç parçası, soldan okunduğunda $t$ ile simge simge uyuşan,
ancak $t$'den daha kısa olan herhangi bir~$s$ dizisidir. Örneğin $[a
\circ {}$, $[a \circ b]$'nin bir öz başlangıç parçasıdır; fakat $[b
\circ {}$ (ikinci simgede uyuşmazlar) ve $[a \circ b]$ (aynı
uzunluktadırlar) değildir.

\begin{prop}\ollabel{prop:initial}
  Bir düzgün~$t$ teriminin her öz başlangıç parçası, $]$ sayısından
  daha fazla sayıda $[$ içerir.
\end{prop}

\begin{proof}
  ~$t$ üzerinde tümevarımla:
  \begin{enumerate}
  \item $t$ tek başına bir harftir: O zaman $t$'nin hiçbir öz başlangıç
    parçası yoktur.
  \item Bazı düzgün $s_1$ ve~$s_2$ terimleri için $t = [s_1 \circ
    s_2]$ olsun. $r$, $t$'nin bir öz başlangıç parçasıysa birkaç olasılık
    vardır:
    \begin{enumerate}
    \item $r$ yalnızca $[$'dir: O zaman $r$, sahip olduğu $]$ sayısından
      bir fazla $[$ içerir.
    \item $r$,~$s_1$'in bir öz başlangıç parçası olan $r_1$ için
      $[r_1$'dir: $s_1$ bir düzgün terim olduğundan, tümevarım hipotezi
      gereği $r_1$, $]$'den daha fazla $[$ içerir; aynı şey $[r_1$ için
      de doğrudur.
    \item $r$, $[s_1$ ya da $[s_1 \circ {}$'dur: Önceki sonuca göre
      ~$s_1$'deki $[$ ve $]$ sayıları eşittir; dolayısıyla $[s_1$ ya da
      $[s_1 \circ {}$ içindeki $[$ sayısı,~$]$ sayısından bir fazladır.
    \item $r$,~$s_2$'nin bir öz başlangıç parçası olan $r_2$ için
      $[s_1 \circ r_2$'dir: Tümevarım hipotezi gereği $r_2$, $]$'den daha
      fazla $[$ içerir. Önceki sonuca göre~$s_1$'deki $[$ ile $]$
      sayıları eşittir. Öyleyse $[s_1 \circ r_2$ içindeki $[$ sayısı,
      ~$]$ sayısından büyüktür.
    \item $r$, $[s_1 \circ s_2$'dir: Önceki sonuca göre $s_1$'deki $[$
      ve $]$ sayıları eşittir; aynı şey~$s_2$ için de geçerlidir. Öyleyse
      $[s_1 \circ s_2$ içinde,~$]$ sayısından bir fazla $[$ vardır.
    \end{enumerate}
  \end{enumerate}
\end{proof}

